\documentclass[10pt,a4paper]{amsart}
\usepackage[margin=19mm]{geometry}
\usepackage{amsmath,amssymb,mathtools}
\usepackage[scr=rsfs,cal=euler]{mathalfa}
\usepackage{xfrac}
\usepackage[unicode,hidelinks,bookmarks=false]{hyperref}
\newcommand{\ie}{i.e.\ }
\newcommand{\cf}{cf.\ }
\newcommand{\resp}{respectively\ }
\newcommand{\Subsection}{\subsection}
\providecommand{\nobreakdash}{\nobreak\hbox{-}}
\setlength{\emergencystretch}{2em}
\multlinegap0pt
\usepackage{amssymb}
\usepackage{dsfont}
\usepackage{natbib}
\newtheorem{prop}{Proposition}
\newtheorem{lem}{Lemma}
\newcommand{\nn}{\nonumber} 
\title{Semiclassical localization of Schrödinger's eigenfunctions}
\author{S\'ebastien Campagne}
\date{Source: arXiv:2602.07128v1 (CC BY 4.0)}
\begin{document}
\maketitle

\noindent\emph{Source and licence.} Semiclassical localization of Schrödinger's eigenfunctions. Version arXiv:2602.07128v1 (CC BY 4.0), distributed by arXiv under the Creative Commons Attribution 4.0 International licence, http://creativecommons.org/licenses/by/4.0/, as recorded in the arXiv licence metadata for this exact version and shown on the abstract page https://arxiv.org/abs/2602.07128v1. Full licence notice: \texttt{licenses/arxiv-2602.07128-CC-BY-4.0.txt}.\par\smallskip

\noindent\textit{Editorial scope.} Counted unit: the article's lem lem (source0.tex lines 1778--1795). The article's complete original proof of it is delivered with it (source0.tex lines 1797--1849, one \texttt{proof} environment of 53 lines). No mathematics is written, repaired or re-derived: the statement and the proof are the article's own lines, in the article's own notation, and no retained line is substituted or re-flowed.  Retained uncounted as the source-local context the proof needs: lines 677--696; lines 709--712.  Nothing was omitted from any retained span, so the package carries no omission record.  No retained line contains a citation, so the wrapper carries no bibliography and no bibliography entry.  No retained line contains a figure environment, an image-inclusion command or drawing code, so no image is delivered and the package declares no asset dependency.  Editorial formatting only, in addition: an AMS \texttt{amsart} preamble that restates the article's own declarations and macros used by the retained lines, namely the environments \texttt{prop}, \texttt{lem} on separate counters, together with the standard packages those lines need.  The retained environments print in this document as Proposition 1, Lemma 1 in order of appearance, whereas the article prints the same items with the numbering of its own full document.  The complete native source of the article as distributed in the arXiv e-print for this exact version is bundled unchanged as \texttt{sources/194152/source0.tex}.
\begin{prop}
\label{prop: def_de_psi_h}
There exists a $K_h$-quasi-conformal homeomorphism $\psi_h$ of the plane (where the dilatation coefficient $K_h$ is given by $K_h = \frac{1 + |\mu_h|}{1 - |\mu_h|}$) satisfying the following properties:
\begin{itemize}
    \item $\psi_h \in W_{loc}^{1,2}$,
    \item $\frac{\partial \psi_h}{\partial \bar{z}} = \mu_h \frac{\partial \psi_h}{\partial z}$,
    \item $K_h \leq 1 + \tilde{C} \epsilon(h)^2$, where $\tilde{C} > 0$,
    \item $\psi_h(0) = 0$ and $\psi_h(\infty) = \infty$.
\end{itemize}

Moreover, the following estimate holds for the derivative of $\psi_h$:
\begin{equation}
\|\partial_z \psi_h - 1\|_p \leq \frac{C_p \|\mu_h\|_p}{1 - C_p \|\mu_h\|_\infty},
\end{equation}
where the constant $C_p$ is defined by
\begin{equation}
C_p := \frac{\pi^2}{4} \cdot \frac{1}{\left(\frac{p}{p-1}\right)^{2/p} - 1} \lesssim_{p \to +\infty} p^2.
\end{equation}
This estimate is valid when $C_p \|\mu_h\|_\infty < 1$, which holds for large $p$ since $\epsilon(h)$ is small enough.
\end{prop}

\begin{equation}
\frac{1}{16} \left| \frac{z_1 - z_2}{R_0} \right|^{K_h} \leq \frac{|g_h(z_1) - g_h(z_2)|}{R_0} \leq 16 \left| \frac{z_1 - z_2}{R_0} \right|^{1/K_h}.
\label{eq: InegdeMori}
\end{equation}

\begin{lem}
		\hfill\\
		Let $K$ be a compact set of $B(0,R_0)$ such that $d(K, B(0,R_0)^c)>\delta>0$ with $\delta>0$. For $h$ small enough:
		\begin{itemize}
			\item There is a constant $\tilde{A}>0$ independent of $h$ such that:
			
			\begin{equation}
				\sup_{z\in \psi_h(K)}|\partial_z R_h(z)|\leq \tilde{A}
			\end{equation}
			
			\item There is a constant $\tilde{B}>0$ independent of $h$ such that:
			
			\begin{equation}
				\frac{1}{\tilde{B}}\leq \inf_{z\in K}|\partial_z R_h(z)|
			\end{equation}
			
		\end{itemize}
	\end{lem}

	\begin{proof}
		
		
		According to proposition \ref{prop: def_de_psi_h} and the Sobolev embedding  $W^{1,3}\subset C^{0,1/3}$, we can show that for $h$ small enough:
		
		\begin{equation}
			||\psi_h- Id||_{L^\infty(B(0,R_0))}\leq R_0^{1/3}||\partial_z\psi_h-1||_{L^3(B(0,R_0))}\longrightarrow 0.
		\end{equation}
		
		So for all $a>0$ small, we have for $h$ small enough, for all $z\in B(0,R_0)$,
		
		\begin{equation}
			\label{eq: eq_sur_psih_pour_Rh}
			|\psi_h(z)-z|\leq a.
		\end{equation} 
		
		\begin{itemize}
			
			\item $R_h$ is a bi-holomorphic mapping from $\psi_h(B(0,R_0))$ to $B(0,R_0)$. According to Cauchy's estimate, for all $\delta>0$, there is a constant $A$ independent of $h$, such that for all compact set $K$ of $B(0,R_0)$ with $d(K, B(0,R_0)^c)>\delta$, we have
			
			\begin{equation}
				\sup_{z\in \psi_h(K)}|\partial_z R_h(z)| \leq A \sup_{z\in \psi_h(B(0,R_0))}|R_h(z)|\leq A R_0
			\end{equation}
			(because $\psi_h$ send $K$ on a compact set $\tilde{K}$ such that  $d(\tilde{K}, \psi_h(B(0,R_0))^c)>\tilde{\delta}=\delta-a>0$ for $h$ small enough).
			
			
			\item Since $R_h$ is bi-holomorphic, $|\partial_z R_h|>0$. So by Cauchy's estimate
			
			\begin{align}
				&\sup_{z\in \psi_h(K)}\left|\frac{1}{\partial_z R_h(z)}\right|\nn\\
				&= \sup_{z\in R_h\circ\psi_h(K)}|\partial_z R_h^{-1}(z)|\leq B \sup_{z\in B(0,R_0)} | R_h^{-1}(z)|,
			\end{align}
			
			because by the Mori inequality \ref{eq: InegdeMori}, $g_h=R_h\circ\psi_h(K)$ send $K$ on $\tilde{K}$ such that\\ $d(\tilde{K}, \psi_h(B(0,R_0))^c)>\tilde{\delta}>0$ for $h$ small enough.\\
			But $ R_h^{-1}(z)$ send $B(0,R_0)$ on $\psi_h(B(0,R_0))$. By the equation \ref{eq: eq_sur_psih_pour_Rh}, we have:
			
			\begin{equation}
				\sup_{z\in B(0,R_0)} | R_h^{-1}(z)|=\sup_{z\in B(0,R_0)} | \psi_h(z)|\leq a+R_0.
			\end{equation}
			So:
			
			\begin{equation}
				\sup_{z\in \psi_h(K)}\left|\frac{1}{\partial_z R_h(z)}\right|\leq \tilde{B},
			\end{equation}
			
			equivalently:
			
			\begin{equation}
				\frac{1}{\tilde{B}}\leq \inf_{z\in \psi_h(K)}|R_h(z)|.
			\end{equation}
			
		\end{itemize}
	\end{proof}

\end{document}
